\documentclass{article}
\usepackage{amsmath,url}
\title{An old carrier's return obligation and its remaining footprint}
\author{}\date{}
\begin{document}\maketitle

\paragraph{Fixed decision interface.}
At one simulated plant cut, the new request needs the left docking cell
and a distinct third carrier should continue. The incumbent delivery
may already be credited while its physical footprint, return obligation,
and authority to re-enter remain. The strict frozen guard requires actual
vacancy, no old in-flight entry, a closed scope for new old entries,
and confirmed discharge of the incumbent duty. The fixture assumes a
selected operator controls the declared scopes and can read the native
grid at this cut; this is not an assumption that a remote coordinator
sees those facts in a deployed installation. A protected sibling return
is evaluated separately from full retirement. It may run at its own
drop cell if it cannot touch the protected dock before the wave ends.

\paragraph{Concrete source effects.}
The exact old warehouse assembly from the parent study uses native
RWARE rewards and positions. The incumbent has a delivery payout at the
left dock and remains loaded there. A local turn and forward step to the
adjacent drop cell produce actual vacancy in the native grid. A toggle
there detaches the shelf and earns the distinct second-half use-and-return
payout. A command to re-enter, issued at the earlier scope, can still
move the old carrier into the left cell after that half is paid.
Similarly, leaving future authority open can allow the old move to be
submitted after the vacancy read. The wrapper represents one ordinary
safety interlock and a queued entrant; RWARE itself is not asserted to
export this administrative API. Its collision rules prevent the new
delivery in the negative cases, not measured trajectories outside the model.

\paragraph{Old duty need not always be fully retired.}
When old access to the protected footprint is closed and all surviving
old action scope fits the disjoint drop cell, the old return can finish
in the same native wave as the new delivery. This is established by
three distinct native half-unit rewards and scoped positions.
No rule treats an unpaid return as paid. The full-retirement guard
refuses that otherwise productive case; the effect-scoped guard
matches the direct fixed-cut permission and result. Apparent disjoint
scope at an earlier read does not suffice if later reassignment is
still possible. In the rebind control, the same initial vacancy and
advertised drop are followed by old entry into the left cell. The
native step blocks the new admitted request and the old drop goes unpaid.
This tests a continued promise of availability by the cut, not a crash
or claim of physical safety on hardware.

\paragraph{Incremental checker on an explicit scoped product.}
The version with genuine closed scope uses the branch where the incumbent
drop and both newer services finish together. The edit enables a
previously dormant reassignment transition; an old possible re-entry
edge reaches the prohibited footprint and then a missed completion.
The paired native position and payout checks supply those concrete
outcomes. The frozen checker from the previous study retains the original
success branch, reconstructs both new rooted counterhistories, and agrees
with the full oracle on this explicit finite product. It does not
extract the registry, discover possible native transitions automatically,
or establish an end-to-end source model. Both the graph interface and the
local authority are declared, rather than inferred from the planner's
display. A new observation-equality or authority model would require its
own compilation and soundness argument.

\paragraph{Comparison.}
This is a development counterpart to the old generation-cut replay.
A stale packet can persist past transmitter withdrawal; similarly
a carrier can still occupy or re-enter a resource after its reported
completion. The concrete carriers of the dependency differ
(callback stages versus footprint and residual return). One strict
abstract state test does not establish a common operational proof for
arbitrary resources. The scoped direct comparator ties.
No independent selection, empirical or field claim, or advance over
an established authority-and-dependency planner is claimed.
\end{document}
